\documentclass[11pt]{article}
\usepackage[margin=0.78in]{geometry}
\usepackage{booktabs}
\usepackage{array}
\usepackage{amsmath}
\usepackage{float}
\usepackage[table]{xcolor}
\usepackage[hidelinks]{hyperref}
\usepackage[numbers,sort&compress]{natbib}
\usepackage{microtype}
\usepackage{seqsplit}

\definecolor{accent}{HTML}{2F6F7E}
\definecolor{light}{HTML}{EAF4F6}
\definecolor{warn}{HTML}{FFF2CC}
\hypersetup{colorlinks=true,linkcolor=accent,citecolor=accent,urlcolor=accent}
\newcolumntype{L}[1]{>{\raggedright\arraybackslash}p{#1}}
\newcolumntype{C}[1]{>{\centering\arraybackslash}p{#1}}
\setlength{\parindent}{0pt}
\setlength{\parskip}{0.46em}

\title{\textbf{Technical Methods Report}\
Released-model reproduction of Nesso-1, Boltz-2, and FlashBind\
on the FEP+ four-kinase benchmark}
\author{Yusheng Cai}
\date{Week of August 17, 2026}

\begin{document}
\maketitle
\begin{center}
  \small Reproducible workflow:
  \href{https://github.com/Yusheng-cai/BindingAffinityPredictor/tree/main/weeks/2026-W34/code}%
       {\texttt{weeks/2026-W34/code}}
\end{center}

\section*{Purpose and scope}

This document records the exact computational procedure I used to reproduce
the released Nesso-1, Boltz-2, and FlashBind affinity results on the public
FEP+ four-kinase benchmark. It is a methods companion to the explanatory Week
34 report. The emphasis here is operational reproducibility: input identity,
model and checkpoint versions, inference settings, native outputs, validation
gates, metric definitions, hardware, and interpretation limits.

The three calculations do not have identical levels of independence:

\begin{itemize}
  \item \textbf{Nesso-1} was run from reconstructed protein sequences and
  stereochemical ligand SMILES. Protein ESM-2 embeddings were generated once
  per unique chain and reused across ligands.
  \item \textbf{Boltz-2} was run from the same reconstructed sequences and
  SMILES, with ColabFold/MMseqs2 MSAs generated once per unique protein chain.
  Boltz-2 generated both protein and ligand coordinates.
  \item \textbf{FlashBind} was run from the authors' released FEP+4 archive.
  That archive already contained experimental PDB protein coordinates,
  FABind+-predicted ligand positions, ESM3 protein representations, and
  TorchDrug ligand representations. I reran the two released affinity
  scorers, but I did not rerun FABind+ docking or feature generation.
\end{itemize}

Consequently, the affinity scores can be compared on the same 87 compound
identities, but the workflows should not be described as equally independent
end-to-end reproductions.

\section{Benchmark and frozen input policy}

I used the historical OpenFF Protein--Ligand Benchmark revision
\texttt{\seqsplit{da7c3372256446222e424368be38ef3d2b55a67b}}, which contains the complete
87 neutral compounds used in the model papers \citep{hahn2022bindingbenchmarks}.
The current project manifest is
\path{data/manifests/fepplus4_87.json}, with SHA256
\texttt{\seqsplit{5f0fb23bb97ea84f3da153e050d5913b97da9bba091fb84cd556b9894fdac1e1}}.

\begin{table}[H]
\centering
\small
\begin{tabular}{L{1.35in}C{0.65in}L{2.55in}L{0.95in}}
\toprule
Target series & $N$ & Protein construct & Experimental endpoint \\
\midrule
CDK2 & 16 & CDK2, 303 aa; cyclin A2, 258 aa & IC50 \\
TYK2 & 16 & TYK2, 302 aa & Ki \\
JNK1 & 21 & JNK1, 370 aa; JIP1 peptide, 11 aa & IC50 \\
p38$\alpha$ & 34 & p38 kinase, 372 aa & IC50 \\
\bottomrule
\end{tabular}
\caption{Frozen affinity benchmark. Each record pairs the indicated target
construct with one ligand.}
\end{table}

Protein sequences were taken from the deposited benchmark complexes. I kept
the cyclin A2 partner for CDK2 and the JIP1 peptide for JNK1. Ligands were
represented by canonical, stereochemically explicit SMILES and restricted to
formal charge zero. Experimental values were converted to $\mu$M and stored as
\[
  y_i=\log_{10}\!\left(\frac{C_i}{\mu\mathrm{M}}\right).
\]
Smaller values correspond to stronger measured potency. The IC50 and Ki labels
were retained as different assay types; neither was relabeled as an
equilibrium binding free energy.

The experimental labels were not supplied to Nesso-1 or Boltz-2 inference.
The FlashBind archive included a label file, but I did not use it to generate
predictions or as the source of the primary analysis labels. I joined all
predictions to the frozen project manifest only after output collection.

\section{Common execution and audit procedure}

All three models were executed on one NVIDIA GeForce RTX 3080 with 10,240 MiB
VRAM and NVIDIA driver 535.230.02. Each model used an isolated environment;
external \texttt{PYTHONPATH} settings were unset except when FlashBind's source
directory was deliberately supplied. Every substantive command was wrapped by
\texttt{scripts/run\_with\_provenance.py}, which recorded the resolved command,
Git revision, software environment, timestamps, wall time, GPU observations,
return code, and failure status in \texttt{run.json}.

The same staged rule was applied to each workflow:
\begin{enumerate}
  \item Verify the pinned source, checkpoint, input, and environment files.
  \item Run a small smoke or target-level pilot and require finite native
  outputs.
  \item Expand to the frozen 87 records without changing scientific settings.
  \item Preserve native model output files before parsing or normalization.
  \item Mark missing or failed records explicitly; do not silently omit them.
  \item Join experimental labels by immutable sample ID only during analysis.
\end{enumerate}

The top-level inference entry point was:
\begin{verbatim}
bash weeks/2026-W34/code/04_run_inference.sh --model=all --execute
\end{verbatim}
The model-specific commands described below are defined in that script and
are printed without execution when \texttt{--check} is used.

\section{Nesso-1 procedure}

\subsection{Pinned implementation and input construction}

I used Nesso-1 version 1.0.0 at source revision
\texttt{\seqsplit{f0156e9a22326448684bae09ee96f73415902dcd}}. The released checkpoint was
resolved to Hugging Face commit
\texttt{\seqsplit{1896c84c7186c506c7efd79051480809d51098bf}}; its model SHA256 was
\texttt{\seqsplit{9928a8a824d147d665e76656804af1cd91c86731516d0b561f9fd1c91ee45622}}.
Protein sequence embeddings used
\texttt{facebook/esm2\_t33\_650M\_UR50D} at revision \texttt{08e4846}.

One YAML input was generated per protein--ligand record. Each contained the
exact protein chain sequence or sequences, one ligand as stereochemical SMILES,
and an affinity request naming that ligand as the binder. No MSA, experimental
protein coordinates, known pocket, template, or experimental label was
provided. ESM-2 embeddings were generated once for each unique protein chain
and reused across every ligand in the same target series.

\subsection{Inference settings and model flow}

The released model was run with seed 42, mixed bfloat16 precision, five
recycling steps, and the default two-stage pocket refinement. The first trunk
pass processed the complete tokenized protein--ligand system and produced
protein--ligand distance distributions. Protein tokens with predicted ligand
distance up to 22~\AA{} were considered for a refinement crop containing at
most 256 total protein-residue and ligand-heavy-atom tokens. Five subsequent
recycled trunk passes refined this cropped pair representation. A tighter
15~\AA{} predicted-distance region was then supplied to the two released
affinity heads \citep{shenoy2026nesso1}.

The resolved settings were:
\begin{table}[H]
\centering
\small
\begin{tabular}{L{2.25in}L{3.85in}}
\toprule
Setting & Value \\
\midrule
Accelerator / devices & GPU / 1 \\
Precision & \texttt{bf16-mixed} \\
Seed & 42 \\
Recycling & 5 feedback steps after the initial pass \\
Pocket refinement & enabled; 22~\AA{} predicted-distance cutoff \\
Refinement token budget & 256 total tokens \\
Final affinity protein cutoff & 15~\AA{} \\
Optional accelerated kernels & disabled with \texttt{--no\_kernels} \\
\bottomrule
\end{tabular}
\caption{Nesso-1 settings used for all four target series.}
\end{table}

For each record, I preserved the ensemble mean
\path{affinity_pred_value}, the component outputs
\path{affinity_pred_value1} and \path{affinity_pred_value2}, the
binary binder probability, and interface entropy diagnostics. The continuous
ensemble mean was used for affinity ranking; binder probability was analyzed
as a separate output and was never substituted for the continuous value.

All 87 records completed. The four target runs required a summed 881.31 s of
wall time; the largest observed total GPU-memory use was 9,292 MiB. These are
local execution measurements, not architecture-independent speed benchmarks.

\section{Boltz-2 procedure}

\subsection{Pinned implementation, inputs, and MSAs}

I used Boltz 2.2.1 source revision
\texttt{\seqsplit{b1ebfc46ecf57f5414e0d1a6f9027bbb122c53bc}}. The structure checkpoint
\path{boltz2_conf.ckpt} had SHA256
\texttt{\seqsplit{090e82ac8c92f5e943fa1b39e7410a44027bea7243c0bbb3caa67a77fc1428e1}};
the affinity checkpoint \path{boltz2_aff.ckpt} had SHA256
\texttt{\seqsplit{dcc5cd3722b1c9eaa34267e4ae32f55cbbf1963f4c19319381ccfa30fdd2ca9e}}.

The sequence and SMILES identities were the same as for Nesso-1. In addition,
I submitted each unique protein chain sequence to the public ColabFold
MMseqs2 server and reused the resulting MSA for every ligand in that target.
No structural template, pocket constraint, experimental coordinate, or
experimental affinity label was supplied \citep{passaro2025boltz2}.

The parsed MSA pool was capped at 8,192 rows. Passing the complete retained MSA
to the trunk exceeded the available 10 GiB GPU memory, so the reported
\texttt{boltz2\_msa1024} protocol enabled Boltz's released stochastic MSA
subsampling option and selected 1,024 rows from that pool for each model pass.
The original full-MSA failure was preserved as a hardware-feasibility result;
it was not counted as an accuracy result.

\subsection{Structure and affinity sampling}

Boltz-2 jointly generated protein and ligand heavy-atom coordinates. For each
record, I used three structure recycles, one structure diffusion sample with
200 sampling steps, step scale 1.5, and one sample in memory at a time. The
affinity calculation used five diffusion samples with 200 sampling steps. I
used seed 42, disabled optional accelerated kernels, did not enable potentials,
and did not apply the optional molecular-weight affinity correction.

\begin{table}[H]
\centering
\small
\begin{tabular}{L{2.25in}L{3.85in}}
\toprule
Setting & Value \\
\midrule
Structure recycles & 3 \\
Structure samples / steps & 1 sample / 200 steps \\
Affinity samples / steps & 5 samples / 200 steps \\
Step scale & 1.5 \\
MSA pool / trunk sample & at most 8,192 parsed / 1,024 sampled rows \\
Templates / constraints / potentials & none / none / disabled \\
Affinity molecular-weight correction & disabled \\
Output coordinates & mmCIF \\
Seed & 42 \\
\bottomrule
\end{tabular}
\caption{Boltz-2 MSA-1024 settings used for all 87 records.}
\end{table}

I retained each sampled complex, confidence outputs, continuous
\path{affinity_pred_value}, and \path{affinity_probability_binary}.
The continuous field was used for affinity ranking; the binder probability
remained separate. All 87 MSA-1024 records completed. The four target runs
required a summed 5,921.84 s, with a maximum observed total GPU-memory use of
9,988 MiB.

\section{FlashBind released-pose scoring procedure}

\subsection{Precisely what was and was not rerun}

I used FlashBind source revision
\texttt{\seqsplit{f161268176237ab6ce5757031a8c1b93937e0d37}}. The two released continuous
value checkpoints came from Hugging Face revision
\texttt{\seqsplit{fa6362c85b3350109af4634e3e2c364644338b3c}}; their SHA256 values were
\texttt{\seqsplit{ca87f84dedd4642da693a508b1119c802caa370e356dad31c47b8a708821e553}}
and
\texttt{\seqsplit{59fa2686b0859a632862aec21ea69f2bebb1ab104e48606fe739c929368d488d}}.

The released FEP+4 archive was pinned at revision
\texttt{\seqsplit{50b1511e080236d80f9b8ff1e4d0cd38bb2480b9}}. For each record it already
provided the inputs needed by the final scorer:
\begin{itemize}
  \item one experimental PDB protein structure shared within each target;
  \item one FABind+-predicted ligand position and orientation in SDF format;
  \item a precomputed ESM3 representation for the protein; and
  \item precomputed TorchDrug atom features for the ligand.
\end{itemize}

I therefore reran only the FlashBind scoring stage. I did not search the PDB,
prepare the protein, rerun FABind+, regenerate ligand coordinates, rerun ESM3,
or regenerate TorchDrug features. This distinction is essential when comparing
the result with Boltz-2, which generated the complete complex in this study
\citep{jiang2026flashbind}.

Before inference, I matched all 87 archive molecules to the frozen manifest by
target and InChIKey. All 87 chemical identities matched. The archived labels
matched the project labels with maximum absolute difference
$4.44\times10^{-16}$, and the four archived PDB files were byte-identical to
the local benchmark copies. These were identity audits; the archived labels
were not model inputs.

\subsection{Scoring and calibration}

FlashBind was run in continuous \texttt{value} mode. The scorer used the
supplied three-dimensional ligand position to construct a local protein region
with the released 20~\AA{} distance-threshold setting. Both value checkpoints
were evaluated and their native outputs were averaged. A one-record CDK2 smoke
test first required finite \texttt{pred\_value}, \texttt{pred\_value\_raw},
and molecular weight; the complete 87-record run began only after this passed.

The raw neural-network output was retained as \texttt{pred\_value\_raw}. The
released implementation applies the molecular-weight calibration
\[
 \texttt{pred\_value}
 =2.177929\,\texttt{pred\_value\_raw}
 -0.218761\,\mathrm{MW}^{0.3}+1.472271.
\]
The calibrated \texttt{pred\_value}, documented as approximately
$\log_{10}(\mathrm{IC50}/\mathrm{Ki}/\mathrm{Kd}\ \mathrm{in}\ \mu M)$,
was used for the primary ranking analysis. Lower values indicate stronger
predicted potency. This mixed-assay score is not a thermodynamic
$\Delta G_{\mathrm{bind}}$.

The upstream environment specified PyTorch 2.7.1 with CUDA 12.6, which was
incompatible with the local NVIDIA driver. I used an isolated PyTorch 2.5.1 /
CUDA 12.1 environment while retaining the authors' code, weights, inputs, and
model settings. The calculation is therefore a numerical reproduction of the
released scorer, not a bitwise environment reproduction.

All 87 records completed in 24.90 s, with observed total GPU-memory use of
3,576 MiB. That runtime excludes the upstream FABind+ docking, ESM3 embedding,
and TorchDrug featurization that had already been performed by the authors.

\section{Output normalization and affinity metrics}

Native output JSON files were preserved beneath \texttt{runs/}. I converted
them to a common CSV schema with
\path{scripts/collect_affinity_outputs.py} for Nesso-1 and Boltz-2 and
\path{scripts/collect_flashbind_outputs.py} for FlashBind. Parsing checked
that the designated continuous output and binder probability were finite and
kept the path to the original native file. No per-target rescaling or
post-hoc sign reversal was applied.

For each target $t$, I calculated
\[
 r_t=\operatorname{Pearson}(y_{t,i},\hat y_{t,i}),
 \qquad
 \bar r=\frac{\sum_t n_t r_t}{\sum_t n_t}.
\]
Thus, the primary result is the compound-count-weighted mean of four
within-target Pearson correlations. It does not correlate compounds measured
in unrelated assays across targets. Spearman $\rho$ and Kendall $\tau$ were
calculated as complementary rank statistics. I estimated uncertainty using
2,000 bootstrap replicates with seed 20260817, resampling compounds with
replacement independently inside each target.

\begin{table}[H]
\centering
\small
\begin{tabular}{L{2.35in}C{1.05in}C{1.05in}C{1.05in}}
\toprule
Model/protocol & Pearson $r$ & Spearman $\rho$ & Kendall $\tau$ \\
\midrule
Nesso-1 & 0.807 & 0.819 & 0.641 \\
Boltz-2 MSA-1024 & 0.627 & 0.666 & 0.497 \\
FlashBind, released FABind+ poses & 0.533 & 0.515 & 0.359 \\
\bottomrule
\end{tabular}
\caption{Compound-count-weighted within-target affinity statistics over all
87 records. The corresponding paper Pearson values were 0.80, 0.66, and 0.53.}
\end{table}

\section{Structure-recovery analysis}

The structure analysis used the 16 records for which the exact
protein--ligand identity had a deposited cocrystal. Nesso-1 was excluded
because it does not produce a final Cartesian complex. Boltz-2 generated both
protein and ligand coordinates. FlashBind used an experimental PDB protein
structure and a ligand position previously predicted by FABind+.

For each pair, I matched protein chains by sequence. The experimental pocket
was defined as residues having any heavy atom within 5~\AA{} of the
crystallographic ligand. I calculated three components:
\begin{enumerate}
  \item global protein C$\alpha$ RMSD after fitting all matched C$\alpha$ atoms;
  \item pocket protein C$\alpha$ RMSD after fitting matched pocket C$\alpha$
  atoms; and
  \item symmetry-corrected ligand heavy-atom RMSD after applying the
  protein-pocket transform, without independently fitting the ligand.
\end{enumerate}

For the ligand calculation, I minimized over complete bond-order- and
chirality-compatible graph mappings $\pi$,
\[
 \operatorname{RMSD}_{\mathrm{pose}}
 =\min_{\pi}
 \sqrt{\frac{1}{N}\sum_i
 \left\|R\mathbf{x}_i+\mathbf{t}
 -\mathbf{x}^{\mathrm{ref}}_{\pi(i)}\right\|^2},
\]
where $R$ and $\mathbf{t}$ were determined only from the protein pocket fit.
This prevents the experimental ligand coordinates from improving the
alignment being evaluated.

\begin{table}[H]
\centering
\small
\begin{tabular}{L{2.65in}C{1.35in}C{1.65in}}
\toprule
Median RMSD over 16 pairs (\AA) & Boltz-2 & FlashBind/FABind+ \\
\midrule
Global protein C$\alpha$, fitted & 1.231 & 0.272 \\
Pocket protein C$\alpha$, fitted & 0.400 & 0.208 \\
Ligand heavy atoms after pocket fit & 0.898 & 0.634 \\
\bottomrule
\end{tabular}
\caption{Protein and ligand coordinate comparison. FlashBind's protein values
compare experimental PDB inputs with ligand-specific cocrystals; they are not
FlashBind protein-prediction errors.}
\end{table}

\section{Reproducible command sequence}

From the repository root, the complete preserved workflow is:
\begin{verbatim}
bash weeks/2026-W34/code/01_setup_models.sh --check
bash weeks/2026-W34/code/02_download_assets.sh --check
bash weeks/2026-W34/code/03_prepare_inputs.sh --check
bash weeks/2026-W34/code/04_run_inference.sh --model=all --check

# Expensive execution, only when assets and GPU are available:
bash weeks/2026-W34/code/04_run_inference.sh --model=all --execute

bash weeks/2026-W34/code/05_collect_outputs.sh --execute
bash weeks/2026-W34/code/06_analyze_affinity.sh --execute
bash weeks/2026-W34/code/07_analyze_poses.sh \
  --execute --allow-rcsb-download
\end{verbatim}

MSA generation on a new machine additionally requires explicit permission to
submit protein sequences to the public ColabFold service:
\begin{verbatim}
bash weeks/2026-W34/code/03_prepare_inputs.sh \
  --execute --allow-msa-server
\end{verbatim}

The scripts preserve completed runs, default to a check or command-preview
mode, and keep checkpoints, MSAs, caches, raw datasets, and raw prediction
directories outside Git. The repository tracks the frozen manifests,
configuration files, adapters, metric implementations, compact tables,
figures, report sources, and tests.

\section{Interpretation limits}

This work establishes that the released checkpoints and preserved inputs can
produce the reported within-target ranking behavior under the documented
protocols. It does not establish universal affinity accuracy, prospective
generalization, or prediction of a standard binding free energy. The
benchmark contains four public congeneric kinase series and may overlap model
training chemistry. IC50 and Ki are assay-dependent quantities. The Boltz-2
result uses a documented 1,024-row MSA trunk sample, and the FlashBind result
uses the authors' already docked FABind+ poses and feature archive.

The technical conclusion is therefore limited but reproducible: Nesso-1 and
Boltz-2 were rerun from reconstructed sequence/SMILES inputs, FlashBind's final
scoring stage was rerun from authors' released prepared complexes, and all
three produced complete 87-record continuous affinity outputs that closely
reproduced their respective paper-level Pearson summaries.

\bibliographystyle{unsrtnat}
\bibliography{references}
\end{document}
